\section{System Model and Problem Formulation}\label{sec:problem}

This section fixes the objects that the rest of the paper reasons about. Table~\ref{tab:notation} collects the notation. Every definition is stated so that it can be checked against the implementation; the module that realises it is named where relevant.

\begin{table}[!t]
\caption{Notation.}
\label{tab:notation}
\centering
\scriptsize
\setlength{\tabcolsep}{3pt}
\begin{tabular}{|L{1.25cm}|L{6.1cm}|}
\hline
Symbol & Meaning \\
\hline
$P,P'$ & Base and variant program versions with a common entry point \\
$c=(P,P')$ & A change pair; $\mathcal{C}$ is a population of change pairs \\
$U_c$ & Changed units: changed lines on either side of $c$ \\
$\mathcal{X}$, $I$ & Declared input domain; executed inputs, a finite multiset over $\mathcal{X}$ \\
$e$ & Execution environment (interpreter, sandbox, limits, platform) \\
$\obs$ & Observation boundary: the channels compared, $\obs\subseteq\{\textit{outcome},\textit{exception},\textit{side\_effect},\textit{api}\}$ \\
$o_{\obs}(P,x;e)$ & Canonical observation of $P$ on $x$ in $e$, restricted to $\obs$ \\
$\Ev$ & Evidence record of one change: per-input records, $U_c$ and metadata \\
$\pi(i)$ & Provenance label of input record $i$ (declared, generated, label-derived) \\
$\iota(\Ev)$ & Integrity flag: $1$ iff the evidence directory validates (Algorithm~\ref{alg:validate}) \\
$\varepsilon(\Ev)$ & Environment-consistency flag \\
$K$ & Controller-held anchor key; never present in an evaluated process \\
$\prof(\Ev)$ & Change-conditioned evidence profile (15 features, Table~\ref{tab:profile}) \\
$\textsf{elig}(\Ev)$ & Eligibility under a blocker policy (v8.0 soft, v8.0s strict, v8.1) \\
$g$, $\Rhat(\Ev)$ & Calibrated risk model; estimated false-assurance risk $g(\prof(\Ev))$ \\
$\tau_\alpha$ & Issuing threshold for target risk $\alpha$ \\
$V(\Ev)$ & Decision procedure; values \DIV{}, \NDO{}, \INS{} \\
$\budget$ & Adjudication budget: resources and procedure for searching for a divergence \\
$A(c;\obs,\budget)$ & Adjudication outcome in $\{\textit{div},\textit{nodiv},\textit{unresolved}\}$ \\
$F$ & False-assurance event, $\mathbf{1}[A=\textit{div}]$ for an issued \NDO{} \\
$\R$, $\cov$ & False-assurance risk and coverage of $V$ on $\mathcal{C}$ \\
$c^\ast$ & Matched coverage at which procedures are compared ($61/258$) \\
$T_j$, $\Delta_j$ & Evidence perturbation $j$ and the estimator's response to it \\
$Y$ & Construction label used in place of $F$ in this study \\
\hline
\end{tabular}
\end{table}

\subsection{Changes, Observations and Divergence}
\begin{definition}[Observable divergence]
Let $c=(P,P')$ be a change, $\obs$ an observation boundary and $e$ an environment. The change \emph{diverges observably on input $x$} if $o_{\obs}(P,x;e)\neq o_{\obs}(P',x;e)$ and the inequality is reproduced on re-execution. An input with this property is a \emph{witness}.
\end{definition}
In \sbg{}, \textit{outcome} is the canonicalised, non-truncated result of the call; \textit{exception} is the typed result kind; \textit{side\_effect} covers captured standard output, files written, module globals and argument mutations; and \textit{api} is the recorded signature of the entry point (\code{sbgv8/adapter.py}).

\begin{definition}[Evidence record]
The evidence for $c$ is $\Ev=(I,U_c,\mathit{meta})$. Each $i\in I$ carries an outcome (\code{SAME}, \code{CONFIRMED\_DIFF}, or a non-comparable status), the executed lines $h_s(i)$ and branch directions on each side $s$, an observation fingerprint $\mathit{fp}(i)$, the captured channels $\mathit{ch}(i)\subseteq\obs$, a determinism flag and repeat count, and a provenance label $\pi(i)$. The metadata contain $\iota(\Ev)$, $\varepsilon(\Ev)$ and the blocker codes that the predecessor engine recorded.
\end{definition}

\begin{definition}[Integrity]
$\iota(\Ev)=1$ if and only if the evidence directory passes validation: exact schemas, binding of the hit sidecar and the run-capture record to the hash-chained log head, sequence contiguity, per-run nonce uniqueness, per-run trust status, and, when a key is configured, an anchor $\mathrm{HMAC}_K$ over the digests of all evidence files (\code{evidence\_schema.validate\_dir}). Integrity is a property of the record, not of the programs.
\end{definition}

\begin{definition}[Eligibility]
Let $\beta(\Ev)$ be the set of blocker codes recorded for $\Ev$ (no clean input, unobservable outcome, nondeterminism, changed region not activated, budget exhausted, environment sensitivity), and let $\Theta$ be the set of codes a policy tolerates. Then $\textsf{elig}(\Ev)$ holds if $\Ev$ contains no witness and $\beta(\Ev)\subseteq\Theta$. The soft policy v8.0 tolerates \{changed region not activated, budget exhausted\}; v8.0s tolerates none; v8.1 additionally requires that the predecessor's fixed gate would issue \NDO{}.
\end{definition}

\subsection{Verdicts and the Decision Procedure}
\begin{definition}[Decision procedure]
A decision procedure is a map $V$ from evidence records to verdicts:
\begin{itemize}
\item \DIVfull{} (\DIV{}): a witness exists;
\item \NDOfull{} (\NDO{}): no divergence was observed and the evidence was judged adequate to report this;
\item \INSfull{} (\INS{}): abstention, because the evidence is untrusted, inadequate, non-reproducible or out of scope.
\end{itemize}
The \sbg{} procedure is
\begin{equation}
V(\Ev)=\begin{cases}
\INS{} & \text{if } \iota(\Ev)=0,\\
\DIV{} & \text{else if a witness exists},\\
\NDO{} & \text{else if } \textsf{elig}(\Ev)\wedge\Rhat(\Ev)\le\tau_\alpha,\\
\INS{} & \text{otherwise}
\end{cases}
\label{eq:V}
\end{equation}
(\code{calibrate.Model.decide}; Algorithm~\ref{alg:decide}).
\end{definition}
A \DIV{} verdict carries its witness and can be checked by replay. An \NDO{} verdict carries no witness, and its reliability is the subject of this paper.

\begin{remark}[No-divergence is not equivalence]\label{rem:equiv}
Program equivalence relative to $\mathcal{X}$ is the statement $\forall x\in\mathcal{X}:\ o(P,x)=o(P',x)$ on \emph{all} channels. An \NDO{} verdict asserts only $\forall i\in I:\ o_{\obs}(P,x_i;e)=o_{\obs}(P',x_i;e)$, for a finite $I$, a restricted $\obs$ and one environment $e$, together with a judgment that the evidence is adequate. Equivalence implies that no divergence can be observed, but the converse fails in three independent ways: an unexecuted input ($\mathcal{X}\setminus I$), an unobserved channel, or an environment-dependent behaviour. \NDO{} is therefore observation-conditioned, and nothing in this paper treats it as \code{PROGRAM\_EQUIVALENCE}. Regression verification~\cite{godlin2009regression,lahiri2012symdiff} addresses the latter question with proofs.
\end{remark}

\paragraph*{Why integrity precedes the witness} The order of the cases in~\eqref{eq:V} is a consequence of the threat model, not a convention. Let an \emph{evidence adversary} be any party that, for any change $c$ and any record content, can produce a record for $c$ with that content, but cannot produce one that passes validation ($\iota=1$) without the key $K$. The program under test, which shares a process with part of the harness, is such a party (Section~\ref{sec:trust}).

\begin{proposition}[Integrity must gate every verdict]\label{prop:integrity}
Let $V$ be any decision procedure and suppose there is a record $\Ev^\ast$ with $\iota(\Ev^\ast)=0$ and $V(\Ev^\ast)=v\neq\INS{}$. Then an evidence adversary can make $V$ return $v$ on every change. If $v=\NDO{}$, the false-assurance risk of $V$ on changes chosen by the adversary is $1$; if $v=\DIV{}$, the adversary can force a divergence verdict on equivalent changes.
\end{proposition}
\emph{Proof:} for any change $c$ the adversary writes a record for $c$ with the content of $\Ev^\ast$. The record fails validation, as $\Ev^\ast$ does, so it is indistinguishable from $\Ev^\ast$ to $V$, and $V$ returns $v$. Choosing only changes that diverge within $\obs$ makes every issued \NDO{} a false assurance; choosing equivalent changes makes every issued \DIV{} spurious.\hfill$\square$

Bounded risk under this adversary therefore requires $V(\Ev)=\INS{}$ whenever $\iota(\Ev)=0$, which is property P8 of Section~\ref{sec:ers}. Because a forged witness is also content, the integrity check must precede the witness check. The price is that a true divergence recorded in tampered evidence yields \INS{} rather than \DIV{}: a denial of a true verdict, never a false assurance. The proposition says nothing about adversaries who hold $K$, or who corrupt the observations before they are recorded; Section~\ref{sec:trust} bounds those cases.

\subsection{The False-Assurance Estimand}
Let $\budget$ be an \emph{adjudication budget}: the resources and procedure an independent adjudicator may use to search for a divergence of $c$ within $\obs$, for example a fixed time allowance, access to the patch and tests, and freedom to construct further inputs. Adjudication returns $A(c;\obs,\budget)\in\{\textit{div},\textit{nodiv},\textit{unresolved}\}$.

\begin{definition}[False assurance, risk and coverage]
For a change on which $V$ issued \NDO{}, the \emph{false-assurance event} is
\begin{equation}
F=\mathbf{1}\left[A(c;\obs,\budget)=\textit{div}\right].
\label{eq:F}
\end{equation}
The \emph{false-assurance risk} and the \emph{coverage} of $V$ on population $\mathcal{C}$ are
\begin{align}
\R&=P\!\left(F=1\,\middle|\,V(\Ev)=\NDO{}\right),\label{eq:R}\\
\cov&=P\!\left(V(\Ev)=\NDO{}\right),\qquad c\sim\mathcal{C}.\label{eq:kappa}
\end{align}
\end{definition}
$\R$ is the selective risk of the \NDO{} decision at coverage $\cov$, in the sense of selective classification~\cite{elyaniv2010foundations,geifman2017selective}. It is defined relative to three declared quantities, $\obs$, $\budget$ and $\mathcal{C}$, and changes when any of them changes. Proposition~\ref{prop:mono} below makes precise in which sense the measured $\R$ is a lower bound.

\paragraph*{Unresolved cases} With $n_d$, $n_u$ and $n_n$ the numbers of issued \NDO{} verdicts adjudicated \textit{div}, \textit{unresolved} and \textit{nodiv}, we report the complete-case estimate $n_d/(n_d+n_n)$, the worst case $(n_d+n_u)/(n_d+n_u+n_n)$ and the optimistic bound $n_d/(n_d+n_u+n_n)$. In this study, unlabelled pairs play the role of unresolved cases.

The dependence of $\R$ on $\obs$ and $\budget$ has a fixed direction under two assumptions. Adjudication is \emph{sound}: it reports \textit{div} only with a reproducible witness. It is also \emph{monotone}: writing $\budget\preceq\budget'$ for a budget at least as generous, if $\obs\subseteq\obs'$ and $\budget\preceq\budget'$, an input that adjudication under $(\obs,\budget)$ finds and that separates the versions under $\obs'$ is also found under $(\obs',\budget')$.

\begin{proposition}[Boundary and budget monotonicity]\label{prop:mono}
Fix the set $S$ of changes on which $V$ issued \NDO{}. If $\obs\subseteq\obs'$ and $\budget\preceq\budget'$, then $\R_{\obs,\budget}(S)\le\R_{\obs',\budget'}(S)$, where $\R_{\obs,\budget}(S)$ is the share of $S$ adjudicated \textit{div} under $(\obs,\budget)$.
\end{proposition}
\emph{Proof:} if a witness $x$ separates $P$ and $P'$ under $\obs$, the observations already differ on a channel in $\obs\subseteq\obs'$, so $x$ also separates them under $\obs'$, and by monotonicity it is found under $(\obs',\budget')$. Hence every $c\in S$ with $A(c;\obs,\budget)=\textit{div}$ also has $A(c;\obs',\budget')=\textit{div}$. The numerator can only grow while the denominator $|S|$ is fixed.\hfill$\square$

A risk reported under a declared boundary and budget is therefore a lower bound on the risk of the same verdicts under any wider boundary or stronger budget. Two consequences follow. Reported risks are meaningless without $(\obs,\budget)$, and risks of procedures that use different boundaries cannot be compared. The converse bound does not hold: nothing observed under $\obs$ bounds the risk under $\obs'$ from above.

\begin{definition}[Matched-coverage risk]
For a scored procedure that issues \NDO{} to eligible records in increasing order of $\Rhat$, $\R(c^\ast)$ is the selective risk when exactly a fraction $c^\ast$ of the population is issued, with fractional issuance among tied scores (\code{stats.matched\_risk}; Algorithm~\ref{alg:matched}). It is defined only if the eligible share is at least $c^\ast$.
\end{definition}

The next result explains why risk must be compared at matched coverage and why random abstention is the right reference.

\begin{proposition}[Random abstention preserves expected risk]\label{prop:random}
Let $V$ issue \NDO{} on a set $S$ of $n$ changes, of which $m$ have $F=1$. Let $V_k$ keep a uniformly random subset of $S$ of fixed size $k\le n$ and abstain on the rest. Then the expected risk among the verdicts kept by $V_k$ is $m/n$, the risk of $V$.
\end{proposition}
\emph{Proof:} each element of $S$ is kept with probability $k/n$, so the expected number of kept false assurances is $mk/n$. Because the denominator $k$ is fixed, the expected risk is $(mk/n)/k=m/n$.\hfill$\square$

Abstention therefore lowers the \emph{number} of false assurances without lowering their \emph{rate} unless it is informative. A procedure that issues fewer verdicts can show fewer false assurances while being no better than chance at choosing which verdicts to withhold. In this paper, B6 (random abstention to $c^\ast$) instantiates $V_k$, and its expected risk equals that of plain differential testing (Section~\ref{sec:results}).

\subsection{Estimating Risk from Evidence}
\sbg{} estimates risk per verdict through a profile map $\prof$ (Section~\ref{sec:approach}) and a calibrated model $g$:
\begin{equation}
\Rhat(\Ev)=g\!\left(\prof(\Ev)\right)\approx P\!\left(F=1\,\middle|\,\prof(\Ev),\ \textsf{elig}(\Ev)\right).
\label{eq:rhat}
\end{equation}
The threshold $\tau_\alpha$ trades coverage against risk. If $g$ were perfectly calibrated, the risk among issued verdicts would equal the mean of $\Rhat$ over them. The evaluation tests whether this holds approximately.

\subsection{Evidence Perturbations}
\begin{definition}[Perturbation response]
A perturbation $T_j$ maps an evidence record to a modified copy. For a risk estimator $\Rhat$ and procedure $V$, the response is
\begin{equation}
\Delta_j(\Ev)=\Rhat\!\left(T_j(\Ev)\right)-\Rhat(\Ev),
\label{eq:delta}
\end{equation}
recorded together with $V(\Ev)$, $V(T_j(\Ev))$ and both profiles. The Evidence Response Signature (\ers{}) of $\Ev$ is the vector $\big(\Delta_j(\Ev),V(T_j(\Ev))\big)_j$. Each $T_j$ has a declared expected response; a trial passes, fails, or is not applicable (N/A) when $T_j$ cannot be applied to $\Ev$.
\end{definition}

\begin{remark}[Conformance is necessary, not sufficient]\label{rem:vacuous}
The procedure $V_\bot\equiv\INS{}$ with a constant risk estimate, which always abstains, passes every property of Section~\ref{sec:ers} as implemented: the monotone properties and P6 accept \INS{}, P8 requires \INS{}, and the invariances see no change. Conformance to \ers{} is therefore a necessary condition for a sensible estimator, not evidence of usefulness. It must be read together with coverage, and passes that rest on abstention (\emph{abstention credit}) are reported separately.
\end{remark}

The converse also fails: accuracy does not imply conformance.

\begin{proposition}[Conformance and accuracy are independent]\label{prop:indep}
Let $\mathcal{D}$ be the distribution of evidence records on which risk, coverage and calibration are measured, and let $T_j$ be a perturbation with $T_j(\Ev)\notin\operatorname{supp}(\mathcal{D})$ for some $\Ev\in\operatorname{supp}(\mathcal{D})$. (i) There is a procedure that passes every \ers{} property and has coverage $0$. (ii) For every estimator $g$ there is an estimator $g'$ that agrees with $g$ on $\operatorname{supp}(\mathcal{D})$, and so has the same risk, coverage, Brier score, log loss and calibration on $\mathcal{D}$, but fails property $j$ on $\Ev$.
\end{proposition}
\emph{Proof:} (i) is Remark~\ref{rem:vacuous}. For (ii), set $g'=g$ on $\operatorname{supp}(\mathcal{D})$ and choose $g'(T_j(\Ev))$ to violate the pass condition of $j$: for an invariance, $|g'(T_j(\Ev))-g(\Ev)|>0.02$; for a monotone property and $V(\Ev)=\NDO{}$, $g'(T_j(\Ev))<g(\Ev)-10^{-6}$, which keeps the verdict \NDO{} when $T_j(\Ev)$ is eligible, because the lower value still meets the threshold. (If $T_j(\Ev)$ is ineligible, every estimator passes, so the trial cannot separate estimators.) Every metric on $\mathcal{D}$ depends on $g'$ only through its values on the support.\hfill$\square$

Perturbed records are typically off the empirical support: no development record contains two appended inputs that reach only unchanged lines. Aggregate metrics are thus structurally blind to the defects that \ers{} tests, and \ers{} is blind to accuracy. The two must be evaluated separately, which is why RQ2 and RQ4 are distinct questions. Section~\ref{sec:res-rq2} reports an instance of~(ii): a defect that changed no prediction on unperturbed data.

\subsection{Labels Used in This Study}
The estimand is defined by independent adjudication. In the development study, $F$ is replaced by a \emph{construction label} $Y$. For synthetic pairs, $Y$ is the generator's record of whether the injected change alters behaviour on the declared domain. For QuixBugs pairs, $Y$ records whether the benchmark's tests distinguish the buggy and fixed versions; negative pairs are behaviour-preserving variants by construction. These labels are not adjudications under $(\obs,\budget)$, and they were visible during design. In the QuixBugs arm the same declared inputs feed both the label and the evidence. Every empirical risk in this paper is therefore a risk with respect to construction labels on development data. A blinded two-adjudicator protocol that produces $A(c;\obs,\budget)$ has been implemented but not executed (Section~\ref{sec:repro}).
